% Chapter 9: Multi-Formation Theory
% Part III: Multi-Formation and Temporal Theory

\chapter{다중 형성 이론}
\label{ch:multi-formation}

\begin{quote}
\textit{``하나의 형성이 안정적으로 존재하는 것은 에너지 최소화의 산물이다. 그러나 여러 형성이 공존하는 것은 운동학적(kinetic) 현상이다.''}
\end{quote}

\vspace{0.5em}
지금까지 우리는 하나의 응집 장 $u_t : X_t \to [0,1]$이 단일 형성을 어떻게 기술하는지를 분석했다. 제II부의 핵심 결과---존재 정리(T1), 상전이(T8-Core), 수렴(T14)---는 모두 단일 형성의 균형 이론에 속한다.

그러나 세계는 하나의 형성만으로 이루어지지 않는다. 책상 위의 컵과 접시, 하늘의 구름들, 시야에 들어오는 여러 사람---이들은 동시에 존재하는 복수의 응집적 형성이다. 이 장에서는 단일 형성에서 복수 형성으로의 확장이 왜 비자명한 도전인지, 그리고 K-장(K-field) 구조가 어떻게 이 도전에 답하는지를 다룬다.


\section{왜 다중 형성이 어려운가}
\label{sec:multi-why-hard}

\subsection{단일 장의 수축 문제}

제8장의 T14(수렴 정리)는 사실 강력한 유일성 결과를 함축한다. 연결 그래프 위의 단일 에너지 범함수 $\mathcal{E}(u)$는 체적 제약 $\Sigma_m$ 위에서 그래디언트 흐름을 따라 수렴하며, 충분히 높은 $\beta/\alpha$ 비율에서 비자명한 최소점이 존재한다.

문제는 이 최소점의 구조이다. 연결 그래프 위의 전역 에너지 최소점은 \textbf{항상 $K=1$}이다---즉, 하나의 연결된 형성이 에너지를 전역적으로 최소화한다. 이는 직관적으로 당연하다: 분리된 두 형성의 체적을 합쳐서 하나의 큰 형성으로 만들면, 경계 에너지가 감소하기 때문이다.

\begin{center}
\fbox{\parbox{0.85\textwidth}{\centering
\textbf{열역학적 사실 ($K^* = 1$ 보편성)}\\[6pt]
모든 연결 그래프 위에서, 전역 에너지 최소점은 $K = 1$ (단일 형성)이다.\\
따라서 $K > 1$ 형성은 에너지 최소화만으로는 설명할 수 없다.
}}
\end{center}

이것이 ``다중 형성이 어려운'' 근본적 이유이다. 에너지 기반의 변분 이론이 단일 형성의 존재는 자연스럽게 보장하지만, 복수 형성의 공존은 원천적으로 보장하지 못한다.

\subsection{네 가지 설계 대안}

이 문제를 해결하기 위한 네 가지 구조적 접근이 고려되었다:

\begin{enumerate}
    \item \textbf{껍질 벗기기(Peeling)}: 형성을 하나 찾은 뒤 제거하고, 나머지에서 다시 찾는 순차적 방법. 문제: 형성 간 상호작용을 포착하지 못하며, 순서 의존적이다.

    \item \textbf{K-장(K-Field)}: $K$개의 결합된 응집 장을 곱 다양체 위에서 동시에 최적화. 문제: 계산 비용이 $K$배 증가하나, 상호작용을 자연스럽게 모형화한다.

    \item \textbf{스펙트럼 $C_t$ 방법}: 공동소속 행렬 $C_t$의 고유벡터를 이용한 분리. 문제: $C_t$가 진단적 역할로 강등된 후(v2.1) 적합하지 않다.

    \item \textbf{국소화(Localization)}: 에너지 밀도의 극소점에서 국소적으로 형성을 탐지. 문제: 겹치는 형성을 처리하기 어렵다.
\end{enumerate}

\textbf{K-장 구조가 선택된 이유}: 형성 간 반발(repulsion), 단체 참여 제약(simplex constraint), 결합 헤시안(joint Hessian) 분석이 모두 자연스럽게 통합되며, 형성의 구성적 유일성(constructive uniqueness)을 보장한다.

% [Figure placeholder]
\begin{figure}[t]
\centering
\fbox{\parbox{0.85\textwidth}{\centering\vspace{5cm}
\textbf{[그림 9.1]} K-장 구조의 개략도.\\
좌: 단일 장 $u$---하나의 형성만 포착. 중: 껍질 벗기기---순차적이나 상호작용 무시.\\
우: K-장---$K$개의 결합된 장 $u^1, \ldots, u^K$이 반발과 단체 제약 하에 동시 최적화.
\vspace{0.5cm}}}
\caption{다중 형성 문제에 대한 세 가지 접근. K-장 구조(우)가 형성 간 상호작용을 가장 자연스럽게 포착한다.}
\label{fig:k-field-architecture}
\end{figure}


\section{K-장 구조}
\label{sec:k-field-architecture}

\subsection{곱 다양체 $\Sigma^K_M$}

$K$개의 형성을 기술하기 위해, $K$개의 응집 장 $u^1, u^2, \ldots, u^K$를 도입한다. 각 장 $u^k : X_t \to [0,1]$은 $k$번째 형성의 응집 분포를 나타내며, 고유한 체적 제약을 만족한다:
\begin{equation}
    \Sigma^K_M = \Sigma_{m_1} \times \Sigma_{m_2} \times \cdots \times \Sigma_{m_K}, \quad \text{여기서 } \Sigma_{m_k} = \Big\{ u^k \in [0,1]^n : \sum_x u^k(x) = m_k \Big\}
    \label{eq:product-manifold}
\end{equation}

체적 $m_k$는 각 형성의 ``크기''를 결정한다. 기본 설정은 동일 분할 $m_k = m/K$이며, 비대칭 분할도 허용된다.

\subsection{단체 참여 제약}

각 지점 $x \in X_t$는 여러 형성에 동시에 참여할 수 있으나, 총 참여도의 합은 1을 초과할 수 없다:
\begin{equation}
    \sum_{k=1}^{K} w_k(x) \leq 1, \quad w_k(x) \geq 0 \quad \forall x \in X_t
    \label{eq:simplex-constraint}
\end{equation}
여기서 $w_k(x)$는 지점 $x$의 $k$번째 형성에 대한 참여 가중치이다. 실제 구현에서는 이를 연성 장벽(soft barrier)으로 처리한다:
\begin{equation}
    \mathcal{E}_{\mathrm{bar}} = \lambda_{\mathrm{bar}} \sum_{x \in X_t} \max\!\Big(0,\, \sum_{k=1}^{K} u^k(x) - 1\Big)^2
    \label{eq:simplex-barrier}
\end{equation}

직관적으로, 단체 제약은 ``한 지점은 한 형성에만 깊이 참여할 수 있다''는 원리를 수학적으로 표현한다. 전이 대역에서의 약한 참여는 여러 형성에 걸칠 수 있으나, 깊은 내부성은 배타적이다.

\subsection{형성별 연산자}

각 형성 $u^k$는 고유한 닫힘 및 구별 연산자를 가진다:
\begin{align}
    \mathrm{Cl}_k(u^k)(x) &= \sigma\!\big(a_{\mathrm{cl}} \cdot (P_t u^k(x) - \tau_{\mathrm{cl}})\big) \label{eq:per-formation-closure} \\
    D_k(x;\, 1 - u^k) &= |P_t u^k(x) - P_t(1 - u^k)(x)| \label{eq:per-formation-distinction}
\end{align}
이는 각 형성이 독립적인 자기-완성과 자기-대비 역학을 가짐을 의미한다.

\subsection{형성 간 반발}

$K > 1$의 핵심 메커니즘은 형성 간 반발(inter-field repulsion)이다:
\begin{equation}
    \mathcal{E}_{\mathrm{rep}} = \lambda_{\mathrm{rep}} \sum_{j < k} \sum_{x \in X_t} u^j(x)\, u^k(x)
    \label{eq:repulsion-energy}
\end{equation}
반발 에너지는 두 형성의 공간적 겹침에 비례한다. 이는 형성들이 서로 다른 영역을 점유하도록 하는 \emph{구조적 배타성}을 구현한다.

\subsection{전체 K-장 에너지}

$K$-장 전체 에너지는 다음과 같다:
\begin{equation}
    \mathcal{E}_K(u^1, \ldots, u^K) = \sum_{k=1}^{K} \mathcal{E}_{\mathrm{self}}(u^k) + \mathcal{E}_{\mathrm{rep}} + \mathcal{E}_{\mathrm{bar}}
    \label{eq:k-field-energy}
\end{equation}
여기서 $\mathcal{E}_{\mathrm{self}}(u^k)$는 각 형성의 독립적 에너지(닫힘 + 분리 + 경계/형태학)이다.

\begin{example}[5$\times$5 격자 위의 두 범프]
\label{ex:two-bumps}
$5 \times 5$ 격자에서 좌상단과 우하단에 각각 가우시안 범프를 배치한다. 단일 장 최적화는 $K=1$ (하나의 거대한 형성)로 수렴하지만, $K=2$ 장 구조에서 반발 $\lambda_{\mathrm{rep}} = 10$을 부여하면 두 개의 분리된 형성이 안정적으로 공존한다. 각 형성은 고유한 핵심(core), 경계 대역(boundary band), 외부(exterior)를 가지며, 반발은 두 형성의 전이 대역 겹침을 효과적으로 억제한다.
\end{example}

\begin{example}[10$\times$10 격자에서의 K-장 최적화: 단계별 추적]
\label{ex:10x10-k-field}
이 예제에서는 $10 \times 10$ 격자 위에서 $K = 2$ 장 최적화의 전 과정을 단계별로 추적한다. 매개변수는 $\alpha = 1.0$, $\beta = 30$, $a_{\mathrm{cl}} = 3.5$, $\lambda_{\mathrm{rep}} = 10$, $\lambda_{\mathrm{bar}} = 100$이다.

\textbf{초기화} (반복 0): 두 장 $u^1, u^2$를 격자의 좌상단 사분면과 우하단 사분면에 각각 가우시안 범프로 초기화한다:
\begin{align}
    u^1_0(x,y) &\propto \exp\!\big(-\tfrac{(x-2.5)^2 + (y-2.5)^2}{2 \cdot 2^2}\big), \quad \sum u^1_0 = m_1 \nonumber \\
    u^2_0(x,y) &\propto \exp\!\big(-\tfrac{(x-7.5)^2 + (y-7.5)^2}{2 \cdot 2^2}\big), \quad \sum u^2_0 = m_2 \nonumber
\end{align}

\textbf{초기 단계} (반복 1--50): 각 장은 독립적으로 닫힘 에너지를 감소시킨다. $\mathrm{Cl}_k(u^k)$가 이웃 지지를 반영하여 핵심 지점의 응집을 강화하고, 외곽 지점의 응집을 약화시킨다. 이 단계에서 두 형성은 서로를 ``보지 못한다''---겹침 $\omega_{12} \approx 0$이므로 반발 효과가 미미하다.

\textbf{경계 형성 단계} (반복 50--200): 이중 우물 에너지 $W(u) = u^2(1-u)^2$에 의해 $u^k$ 값이 0 또는 1 근방으로 분극된다. 경계 대역이 명확해지며, $\mathsf{Bind}$ 값이 0.5 → 0.85로 빠르게 상승한다.

\textbf{안정화 단계} (반복 200--500): 에너지 감소율이 둔화되며, 두 형성이 안정적인 극소점에 근접한다. 최종 상태에서:
\begin{itemize}
    \item 형성 1: $\mathrm{Core}(u^1) = \{(1,1), (1,2), (2,1), (2,2), (2,3), (3,2)\}$ (6 지점)
    \item 형성 2: $\mathrm{Core}(u^2) = \{(7,7), (7,8), (8,7), (8,8), (8,9), (9,8)\}$ (6 지점)
    \item 형성 간 최소 거리: $d_{\min} = 5$ (그래프 거리)
    \item $\Lambda_{\mathrm{coupling}} = 0.003$ (잘 분리된 체제)
\end{itemize}
\end{example}

% [Figure placeholder]
\begin{figure}[t]
\centering
\fbox{\parbox{0.85\textwidth}{\centering\vspace{6cm}
\textbf{[그림 9.1b]} 10$\times$10 격자에서의 K-장 최적화 과정.\\
4단계 열: 초기화(반복 0), 초기 단계(반복 50), 경계 형성(반복 200), 안정화(반복 500).\\
상단: $u^1$ 히트맵. 하단: $u^2$ 히트맵. 두 형성의 핵심이 점차 강화되는 과정.
\vspace{0.5cm}}}
\caption{10$\times$10 격자에서의 $K=2$ 장 최적화 과정. 초기의 넓은 가우시안 범프가 점차 날카로운 형성으로 수렴한다.}
\label{fig:10x10-optimization}
\end{figure}

\subsection{K-장 에너지의 기울기}

최적화 구현을 위해, $\mathcal{E}_K$의 기울기를 명시적으로 도출한다. 각 형성 $u^k$에 대한 기울기는 세 기여의 합이다:
\begin{equation}
    \nabla_{u^k} \mathcal{E}_K = \underbrace{\nabla_{u^k} \mathcal{E}_{\mathrm{self}}(u^k)}_{\text{독립적 기울기}} + \underbrace{\lambda_{\mathrm{rep}} \sum_{j \neq k} u^j}_{\text{반발 기울기}} + \underbrace{2\lambda_{\mathrm{bar}} \sum_x \max\!\big(0, \sum_l u^l(x) - 1\big) \cdot \mathbf{e}_x}_{\text{장벽 기울기}}
    \label{eq:k-field-gradient}
\end{equation}

독립적 기울기 $\nabla_{u^k} \mathcal{E}_{\mathrm{self}}$는 제6장의 단일 형성 기울기와 동일하다. 반발 기울기는 단순히 다른 형성들의 합이며, 장벽 기울기는 단체 위반이 발생하는 지점에서만 비영이다. 이 세 기여의 독립적 구조가 K-장 최적화의 계산적 효율성의 근거이다.


\section{결합 한계 보조정리}
\label{sec:coupling-bound}

\subsection{결합 헤시안의 구조}

K-장 에너지의 2차 변분(헤시안)은 블록 구조를 가진다:
\begin{equation}
    H_K = \begin{pmatrix}
        H_1 & R_{12} & \cdots & R_{1K} \\
        R_{12}^T & H_2 & \cdots & R_{2K} \\
        \vdots & & \ddots & \vdots \\
        R_{1K}^T & R_{2K}^T & \cdots & H_K
    \end{pmatrix}
    \label{eq:joint-hessian}
\end{equation}
여기서 $H_k$는 $k$번째 형성의 독립적 헤시안이고, $R_{jk}$는 반발에서 오는 교차 항(cross-term)이다. 반발 교차 항은 특히 단순한 형태를 가진다:
\begin{equation}
    R_{jk} = \lambda_{\mathrm{rep}} \cdot \mathrm{diag}(\mathbf{1})
    \label{eq:cross-term}
\end{equation}
즉, 반발의 교차 효과는 모든 지점에서 균일한 상수이다.

\subsection{스펙트럼-반발 양립성 (SR)}

\textbf{핵심 관찰}: 형성의 핵심 지점(core site)에서 반발의 교차 효과는 사실상 무시할 수 있다. 이유는 간단하다---잘 분리된 형성들의 핵심은 공간적으로 떨어져 있으므로, $u^j(x) \approx 0$ for $x \in \mathrm{Core}(u^k)$ ($j \neq k$)이다.

\begin{definition}[스펙트럼-반발 양립성 (SR)]
\label{def:spectral-repulsion}
K-장 구성 $(u^1, \ldots, u^K)$이 스펙트럼-반발 양립적이라 함은, 각 형성의 핵심 지점에서 교차 항이 소멸하는 것이다:
\begin{equation}
    u^j(x) \cdot u^k(x) < \varepsilon_{\mathrm{SR}}, \quad \forall x \in \mathrm{Core}(u^k),\ j \neq k
    \label{eq:sr-condition}
\end{equation}
\end{definition}

\subsection{Weyl 한계}

결합 헤시안의 고유값에 대한 Weyl 부등식은 다음을 제공한다:
\begin{equation}
    \sigma_{\max}(H_K) \leq K \cdot \lambda_{\mathrm{bd}} + \lambda_{\mathrm{rep}}
    \label{eq:weyl-bound}
\end{equation}
이는 결합 헤시안의 스펙트럼 범위가 형성 수 $K$에 선형적으로 증가함을 의미한다. SR 조건 하에서는 더 강한 결과를 얻는다:

\begin{theorem}[결합 한계 보조정리]
\label{thm:coupling-bound}
SR 조건을 만족하는 K-장 구성에서, 각 형성 $u^k$의 에너지 기울기는 핵심으로부터의 거리에 따라 지수적으로 감쇄한다:
\begin{equation}
    |\nabla_{u^k} \mathcal{E}_K(x)| \leq C \cdot e^{-\gamma \cdot d(x, \mathrm{Core}(u^k))}
    \label{eq:gradient-decay}
\end{equation}
여기서 $C, \gamma > 0$은 $\lambda_{\mathrm{rep}}, \beta, \alpha$에 의존하는 상수이다.
\end{theorem}

직관적 의미: 각 형성은 자신의 핵심 근처에서만 ``활동적''이며, 멀리 떨어진 영역에서의 영향은 지수적으로 감소한다. 이것이 잘 분리된 형성들이 사실상 독립적으로 행동하는 수학적 근거이다.

\begin{proof}[결합 한계 보조정리의 증명 스케치]
$u^k$의 에너지 기울기를 분해하면:
\begin{equation}
    \nabla_{u^k} \mathcal{E}_K(x) = \nabla_{u^k} \mathcal{E}_{\mathrm{self}}(x) + \lambda_{\mathrm{rep}} \sum_{j \neq k} u^j(x) \nonumber
\end{equation}
핵심에서 멀어진 지점 $x$에서, $u^k(x) \approx 0$이므로 $\nabla_{u^k} \mathcal{E}_{\mathrm{self}}(x) \approx 0$이다. 반발 항 $\sum_{j \neq k} u^j(x)$도, SR 조건에 의해 핵심 $\mathrm{Core}(u^j)$에서 멀어지면 지수적으로 감쇄한다 (이는 개별 형성의 경계 감쇄 $u^j(x) \leq C' e^{-\gamma' d(x, \mathrm{Core}(u^j))}$에서 따른다). 두 기여를 합하면 전체 기울기의 지수 감쇄를 얻는다.

감쇄 상수 $\gamma$는 이중 우물의 볼록 영역 폭과 그래프의 국소적 연결도에 의존한다. $\beta$가 커질수록 $\gamma$도 커진다---경계가 더 날카로울수록 기울기의 국소화가 강해진다.
\end{proof}

\begin{example}[기울기 감쇄의 수치적 검증]
\label{ex:gradient-decay-numerical}
$20 \times 20$ 격자에서 $K = 2$ 형성을 최적화한 후, 형성 1의 핵심($(5,5)$ 근방)으로부터의 거리에 따른 $|\nabla_{u^1} \mathcal{E}_K|$를 측정한다:

\begin{center}
\begin{tabular}{ccc}
\hline
\textbf{거리 $d$} & $|\nabla_{u^1} \mathcal{E}_K|$ & \textbf{이론적 한계} \\
\hline
0 & $1.24 \times 10^{-1}$ & --- \\
1 & $8.7 \times 10^{-2}$ & $C e^{-\gamma}$ \\
2 & $3.1 \times 10^{-2}$ & $C e^{-2\gamma}$ \\
3 & $8.5 \times 10^{-3}$ & $C e^{-3\gamma}$ \\
5 & $4.2 \times 10^{-4}$ & $C e^{-5\gamma}$ \\
8 & $1.1 \times 10^{-6}$ & $C e^{-8\gamma}$ \\
\hline
\end{tabular}
\end{center}
피팅 결과 $\gamma \approx 1.45$로, 이론적 예측과 잘 일치한다.
\end{example}

% [Figure placeholder]
\begin{figure}[t]
\centering
\fbox{\parbox{0.85\textwidth}{\centering\vspace{5cm}
\textbf{[그림 9.2]} 결합 헤시안의 블록 구조와 기울기 감쇄.\\
좌: 두 형성의 결합 헤시안---대각 블록(형성별 헤시안)과 교차 블록(반발).\\
우: 핵심으로부터의 거리에 따른 기울기 크기의 지수적 감쇄.
\vspace{0.5cm}}}
\caption{SR 조건 하에서 결합 헤시안의 교차 항은 핵심 지점에서 소멸하며, 기울기는 핵심으로부터 지수적으로 감쇄한다.}
\label{fig:gradient-decay}
\end{figure}


\section{$\Lambda_{\mathrm{coupling}}$에 의한 체제 분류}
\label{sec:regime-classification}

\subsection{결합 매개변수의 정의}

다중 형성의 행동을 지배하는 단일 무차원 매개변수를 도입한다:
\begin{equation}
    \Lambda_{\mathrm{coupling}} = \max_{j \neq k} \frac{\lambda_{\mathrm{rep}} \cdot \omega_{jk}}{\min(\mu_j, \mu_k)}
    \label{eq:lambda-coupling}
\end{equation}
여기서:
\begin{itemize}
    \item $\omega_{jk} = \langle u^j, u^k \rangle / \min(\|u^j\|^2, \|u^k\|^2)$는 부드러운 겹침(soft overlap)
    \item $\mu_k$는 $k$번째 형성의 스펙트럼 간극(spectral gap)---형성의 안정성 척도
\end{itemize}

$\Lambda_{\mathrm{coupling}}$의 물리적 의미: 형성 간 반발 결합의 강도를 개별 형성의 안정성에 대비하여 측정한 것이다. $\Lambda$가 작으면 반발이 안정성에 비해 약하므로 형성은 독립적으로 행동하고, $\Lambda$가 크면 반발이 안정성을 압도하여 형성 간 상호작용이 지배적이 된다.

\subsection{세 가지 체제}

$\Lambda_{\mathrm{coupling}}$의 값에 따라 세 가지 체제가 구분된다:

\begin{enumerate}
    \item \textbf{잘 분리된 체제} (Well-Separated, $\Lambda < 0.01$): 형성들의 지지(support)가 본질적으로 겹치지 않는다. $\omega_{jk} \approx 0$이므로 $\Lambda \approx 0$. 각 형성은 독립적인 단일 형성처럼 행동하며, 지속성(persistence) 분석이 가장 단순하다.

    \item \textbf{약하게 상호작용하는 체제} (Weakly-Interacting, $0.01 \leq \Lambda < 1/(K-1)$): 형성의 경계 대역이 겹칠 수 있으나, 핵심은 분리되어 있다. 반발의 효과가 감지되지만 형성의 구조적 안정성을 파괴하지는 않는다.

    \item \textbf{강하게 상호작용하는 체제} (Strongly-Interacting, $\Lambda \geq 1/(K-1)$): 형성의 부피(bulk)가 겹치며, 운동학적 장벽(kinetic barrier)을 넘나드는 전이가 가능하다. 형성 병합 또는 소멸이 빈번히 일어난다.
\end{enumerate}

임계값 $\Lambda = 1/(K-1)$은 $K$개 형성의 등분할에서 반발 결합이 개별 형성의 안정성을 압도하는 경계를 나타낸다.

\begin{example}[$\Lambda_{\mathrm{coupling}}$ 계산 실례]
\label{ex:lambda-coupling-calc}
$15 \times 15$ 격자에서 $K = 3$ 형성을 고려하자. 세 형성의 핵심은 각각 $(3,3)$, $(12,3)$, $(7,12)$에 위치한다.

\textbf{단계 1: 부드러운 겹침 계산.} 형성 쌍별 겹침:
\begin{align}
    \omega_{12} &= \langle u^1, u^2 \rangle / \min(\|u^1\|^2, \|u^2\|^2) = 0.0023 \nonumber \\
    \omega_{13} &= \langle u^1, u^3 \rangle / \min(\|u^1\|^2, \|u^3\|^2) = 0.0041 \nonumber \\
    \omega_{23} &= \langle u^2, u^3 \rangle / \min(\|u^2\|^2, \|u^3\|^2) = 0.0038 \nonumber
\end{align}

\textbf{단계 2: 스펙트럼 간극.} 각 형성의 헤시안 최소 고유값:
\begin{equation}
    \mu_1 = 2.31, \quad \mu_2 = 2.18, \quad \mu_3 = 2.45 \nonumber
\end{equation}

\textbf{단계 3: $\Lambda_{\mathrm{coupling}}$ 계산.} $\lambda_{\mathrm{rep}} = 10$일 때:
\begin{equation}
    \Lambda_{\mathrm{coupling}} = \max\!\Big\{\frac{10 \cdot 0.0023}{\min(2.31, 2.18)},\; \frac{10 \cdot 0.0041}{\min(2.31, 2.45)},\; \frac{10 \cdot 0.0038}{\min(2.18, 2.45)}\Big\} = \frac{0.041}{2.18} = 0.0188 \nonumber
\end{equation}

\textbf{판정}: $0.01 \leq 0.0188 < 1/(3-1) = 0.5$ → \textbf{약하게 상호작용하는 체제}.

이 판정은 기하학적 관찰과 일치한다: 형성 1과 3의 경계 대역이 약간 겹치지만(그래프 거리 $\approx 8$), 핵심은 완전히 분리되어 있다.
\end{example}

\subsection{체제 전이의 역학}

$\Lambda_{\mathrm{coupling}}$은 정적 매개변수가 아니다---형성이 이동하거나 변형됨에 따라 시간적으로 변화한다. 특히 두 형성이 접근할 때, $\omega_{jk}$가 증가하여 $\Lambda$가 상승하며, 체제 전이가 일어날 수 있다:
\begin{equation}
    \text{잘 분리된} \xrightarrow{\omega_{jk} \uparrow} \text{약하게 상호작용} \xrightarrow{\omega_{jk} \uparrow\uparrow} \text{강하게 상호작용}
    \label{eq:regime-transition}
\end{equation}

강하게 상호작용하는 체제에 진입하면 형성 병합이 발생할 수 있다. 이 전이의 날카로움은 $\beta$에 의존한다: $\beta$가 클수록 전이가 급격하며(준안정 장벽이 높으므로 한번 넘으면 빠르게 병합), $\beta$가 작으면 전이가 완만하다.

\begin{center}
\fbox{\parbox{0.85\textwidth}{\centering
\textbf{경험적 검증}: 실험 exp46--47에서 69개 구성에 대해 $\Lambda_{\mathrm{coupling}}$ 기반 체제 분류와\\
기하학적 기준(지지 겹침, 핵심 분리)의 일치율이 \textbf{100\%}로 확인되었다.
}}
\end{center}

% [Figure placeholder]
\begin{figure}[t]
\centering
\fbox{\parbox{0.85\textwidth}{\centering\vspace{5cm}
\textbf{[그림 9.3]} $\Lambda_{\mathrm{coupling}}$ 체제 위상 다이어그램.\\
가로축: $\lambda_{\mathrm{rep}}$ (반발 강도). 세로축: $d_{\min}$ (형성 간 최소 거리).\\
세 체제가 색상으로 구분: 잘 분리된(녹색), 약하게 상호작용(황색), 강하게 상호작용(적색).
\vspace{0.5cm}}}
\caption{$\Lambda_{\mathrm{coupling}}$에 의한 체제 분류. 69개 실험 구성이 기하학적 기준과 100\% 일치한다.}
\label{fig:regime-phase-diagram}
\end{figure}


\section{다중 형성은 운동학적이다, 열역학적이 아니다}
\label{sec:kinetic-not-thermodynamic}

이 절은 이 장의 핵심 메시지를 담고 있다. 단일 형성 이론(제II부)은 에너지 최소화의 관점에서---즉 \textbf{열역학적(thermodynamic)} 관점에서---완전하다. 그러나 다중 형성 이론은 본질적으로 \textbf{운동학적(kinetic)}이다.

\subsection{$K^* = 1$ 보편성}

\begin{theorem}[$K^* = 1$ 보편성]
\label{thm:k-star-one}
임의의 연결 그래프 $(X, W)$와 임의의 매개변수 $(\alpha, \beta, \lambda_{\mathrm{cl}}, \lambda_{\mathrm{sep}})$에 대해, 전역 에너지 최소점은 $K = 1$이다. 즉, $K > 1$ 형성은 전역 최소점이 아닌 \textbf{준안정 극소점(metastable local minimum)}이다.
\end{theorem}

\begin{proof}[증명의 핵심 아이디어]
$K = 2$인 두 분리된 형성 $u^1, u^2$에서, $u = u^1 + u^2$ (겹침이 없을 때)로 합치면 반발 에너지 $\mathcal{E}_{\mathrm{rep}}$가 0으로 떨어지며, 경계 에너지도 감소한다(경계선의 총 길이가 줄어드므로). 따라서 합쳐진 형성이 더 낮은 에너지를 가진다. 연결 그래프의 경우, 이 논증이 임의의 $K$에 대해 귀납적으로 확장된다.
\end{proof}

이 결과의 함의는 심대하다: \emph{에너지만으로는 다중 형성을 설명할 수 없다.}

\begin{remark}[인지적 유비]
일상 경험에서 이 결과는 직관적으로 이해된다. 책상 위의 컵과 접시는 에너지적으로 ``합쳐져야'' 할 이유가 없다---이들은 각각의 공간적 위치에서 안정적으로 유지된다. SCC의 수학적 언어로, 컵과 접시는 운동학적 장벽에 의해 보호되는 준안정 극소점이다. 에너지가 ``하나의 큰 형성''을 선호하지만, 장벽이 너무 높아서 자발적 병합이 일어나지 않는다.
\end{remark}

\subsection{세 가지 운동학적 기둥}

$K > 1$ 형성의 존재와 안정성을 설명하는 세 가지 운동학적 메커니즘이 있다:

\textbf{(1) 핵생성(Nucleation).} 형성은 균일 상태로부터의 분기(bifurcation)를 통해 탄생한다. T8-Core에 의하면, 임계 비율 $\beta/\alpha > 4\lambda_2/|W''(c)|$를 넘으면 균일 상태가 불안정해지며, 불안정한 방향은 그래프 라플라시안의 \textbf{Fiedler 고유벡터}가 결정한다. $K > 1$ 형성은 더 높은 고유벡터들이 동시에 불안정해질 때 핵생성된다.

\textbf{(2) 준안정성(Metastability).} 일단 형성된 $K > 1$ 형성은 운동학적 장벽에 의해 보호된다. 두 형성이 병합되려면 장벽을 넘어야 하며, 이 장벽의 높이는 $\beta$에 따라 다음과 같이 증가한다:
\begin{equation}
    \Delta E_{\mathrm{barrier}} \propto \beta^{0.89}
    \label{eq:barrier-scaling}
\end{equation}
이 스케일링은 exp55에서 경험적으로 확인되었다. 지수 $0.89$는 고전적 Allen-Cahn 장벽 스케일링 $\beta^1$보다 약간 낮은데, 이는 닫힘 연산자의 비선형 효과 때문이다.

\textbf{(3) 조대화(Coarsening).} 잡음 존재 하에서 $K > 1$ 형성은 궁극적으로 병합하여 $K-1, K-2, \ldots, 1$로 감소한다. 고전적 Allen-Cahn 조대화에서 성장 지수(coarsening exponent)는 $\alpha_c = 1/2$이다: 평균 영역 크기가 $\langle R \rangle \sim t^{1/2}$로 성장한다. SCC에서는 닫힘 연산자의 자기-강화 효과에 의해 조대화가 \textbf{더 느리다}: $\alpha_c < 1/2$ (운동학적 예측 MK-2).

그러나 잘 분리된 형성 쌍은 놀랍도록 안정적이다: exp55에서 5000 반복, 잡음 진폭 $\sigma = 0.01$ 조건에서 병합이 \textbf{한 번도 발생하지 않았다.} 이는 장벽 높이가 잡음 에너지 $\sim \sigma^2 n$을 크게 초과하기 때문이다.

\begin{example}[잡음 안정성 실험]
\label{ex:noise-stability}
$20 \times 20$ 격자, $K = 2$, $\beta = 30$, 형성 간 거리 $d = 8$. 매 반복마다 독립 가우시안 잡음 $\xi \sim \mathcal{N}(0, \sigma^2)$를 각 $u^k(x)$에 주입한 후 $\Sigma_m$으로 재사영한다.
\begin{itemize}
    \item $\sigma = 0.001$: 5000 반복, 병합 0회. 진단 벡터 변동 $< 1\%$.
    \item $\sigma = 0.01$: 5000 반복, 병합 0회. 진단 벡터 변동 $< 5\%$.
    \item $\sigma = 0.05$: 5000 반복, 병합 3회 (반복 2100, 3400, 4800). 병합 시 $\Lambda_{\mathrm{coupling}}$이 0.01을 넘어 강하게 상호작용 체제로 전이됨을 확인.
    \item $\sigma = 0.1$: 100 반복 내에 병합. 장벽 완전 파괴.
\end{itemize}
\end{example}

\begin{center}
\fbox{\parbox{0.85\textwidth}{\centering
\textbf{핵심 통찰}\\[6pt]
다중 형성의 존재는 에너지의 전역 최소화가 아니라 \emph{운동학적 장벽}에 의해 보호된다.\\
$K > 1$ 형성은 준안정적 극소점이며, 장벽 높이는 $\beta^{0.89}$에 비례한다.
}}
\end{center}

\subsection{닫힘 연산자의 준안정성 강화}

운동학적 관점에서 닫힘 연산자 $\mathrm{Cl}_t$의 역할은 특히 중요하다. 닫힘은 형성의 내부 응집을 강화하여 준안정 장벽을 \emph{더 높게} 만든다. 이는 예측 P3(강화된 체류, enhanced dwell)의 다중 형성 버전이다:

순수 Allen-Cahn 에너지($\lambda_{\mathrm{cl}} = \lambda_{\mathrm{sep}} = 0$)와 비교할 때, SCC의 닫힘 항은 장벽을 유의미하게 증가시킨다. 이는 닫힘이 형성의 핵심을 더 ``단단하게'' 만들어, 잡음에 의한 장벽 넘기를 어렵게 하기 때문이다.

% [Figure placeholder]
\begin{figure}[t]
\centering
\fbox{\parbox{0.85\textwidth}{\centering\vspace{5cm}
\textbf{[그림 9.4]} 준안정 장벽 높이의 $\beta$ 스케일링.\\
가로축: $\beta$ (분리 강도). 세로축: $\Delta E_{\mathrm{barrier}}$.\\
실험 데이터(점)와 $\beta^{0.89}$ 피팅(실선). SCC(파란)와 순수 Allen-Cahn(회색) 비교.
\vspace{0.5cm}}}
\caption{장벽 높이의 $\beta$ 스케일링. SCC의 닫힘 항은 순수 Allen-Cahn 대비 장벽을 유의미하게 강화한다.}
\label{fig:barrier-scaling}
\end{figure}


\section{T-Beyond-Weyl 스펙트럼 한계}
\label{sec:beyond-weyl}

\subsection{표준 Weyl 한계의 한계}

그래프 라플라시안의 고유값에 대한 표준 Weyl 한계는 다음을 제공한다:
\begin{equation}
    \sigma_i(L) \leq 2\sqrt{d_{\max}}, \quad \text{($d$-정칙 그래프의 경우)}
    \label{eq:standard-weyl}
\end{equation}
이 한계는 형성의 공존 가능 영역(coexistence window)을 결정하는 데 핵심적이다: $K > 1$ 형성이 공존하려면 충분히 많은 불안정 모드가 필요하며, 이는 $\beta/\alpha$가 연속적인 고유값들의 문턱을 넘어야 함을 뜻한다. Weyl 한계가 느슨하면 공존 영역이 과소평가된다.

\subsection{경계 모드 우위 (BMD)}

SCC 형성의 핵심적 특성은 \textbf{부드러운 국소화(soft localization)}이다. 형성의 핵심 근방에서 에너지 기울기는 급격히 변하며, 이는 헤시안의 고유벡터가 핵심 근방에 국소화됨을 의미한다.

\begin{definition}[경계 모드 우위, BMD]
\label{def:bmd}
형성 $u^k$에 대해, 헤시안 $H_k$의 지배적 고유모드가 $\mathrm{Bd}(u^k)$ 근방에 질량의 대부분을 집중시킬 때, $u^k$는 경계 모드 우위를 가진다고 한다.
\end{definition}

BMD를 이용하면 표준 Weyl 한계를 상당히 개선할 수 있다:

\begin{theorem}[T-Beyond-Weyl]
\label{thm:beyond-weyl}
부드러운 국소화를 이용하면, 특정 그래프 구조에서 유효 고유값 한계가 표준 Weyl 한계 대비 최대 \textbf{33배} 더 조밀해진다. 이는 다중 형성의 공존 영역을 크게 확장한다.
\end{theorem}

직관적으로, Weyl 한계는 그래프 전체의 ``최악의 경우''를 반영하지만, 형성의 부드러운 국소화는 에너지 곡면의 유효 자유도를 줄여서 더 좁은 스펙트럼 범위 내에 더 많은 불안정 모드를 허용한다.

% [Figure placeholder]
\begin{figure}[t]
\centering
\fbox{\parbox{0.85\textwidth}{\centering\vspace{4cm}
\textbf{[그림 9.5]} 스펙트럼 한계 비교.\\
표준 Weyl 한계(점선)와 T-Beyond-Weyl 한계(실선).\\
부드러운 국소화에 의해 유효 고유값 간격이 33배 더 조밀해짐.
\vspace{0.5cm}}}
\caption{T-Beyond-Weyl 한계. 부드러운 국소화는 표준 Weyl 한계를 크게 개선하여 다중 형성의 공존 영역을 확장한다.}
\label{fig:spectral-bound}
\end{figure}


\section{임계 형성 간 거리}
\label{sec:critical-distance}

\subsection{경험적 공식}

두 형성이 안정적으로 공존하기 위한 최소 거리 $d_{\min}^*$는 세 매개변수에 의존한다:
\begin{equation}
    d_{\min}^* = f(\beta,\, a_{\mathrm{cl}},\, \lambda_{\mathrm{rep}})
    \label{eq:d-min-formula}
\end{equation}

일반적으로, $d_{\min}^*$는 다음과 같은 정성적 행동을 보인다:
\begin{itemize}
    \item $\beta$ 증가: $d_{\min}^*$ 감소 (더 날카로운 경계 → 더 좁은 전이 대역 → 더 가까이 공존 가능)
    \item $\lambda_{\mathrm{rep}}$ 증가: $d_{\min}^*$ 증가 (더 강한 반발 → 더 넓은 분리 필요)
    \item $a_{\mathrm{cl}}$ 증가: $d_{\min}^*$ \textbf{감소} (더 강한 닫힘 → 핵심 강화 → 더 가까이 공존 가능)
\end{itemize}

\subsection{닫힘이 $d_{\min}^*$를 30\% 감소시킨다}

가장 주목할 만한 결과는 닫힘의 역할이다:

\begin{center}
\fbox{\parbox{0.85\textwidth}{\centering
\textbf{CN14: 닫힘은 다중 형성 안정성을 확장한다}\\[6pt]
닫힘 연산자는 임계 형성 간 거리 $d_{\min}^*$를 약 \textbf{30\%} 감소시킨다.\\
예: $\beta = 30$에서 순수 Allen-Cahn은 $d \approx 7$을 요구하나,\\
SCC는 $d \approx 5$에서도 두 형성이 안정적으로 공존한다.
}}
\end{center}

이 결과는 닫힘 연산자가 단일 형성의 안정성(준안정 장벽 강화)뿐 아니라, 다중 형성의 \emph{공존 영역}도 확장한다는 것을 의미한다. 닫힘은 각 형성의 핵심을 더 단단하게 만들어, 가까운 거리에서도 구별된 핵심을 유지할 수 있게 한다.

% [Figure placeholder]
\begin{figure}[t]
\centering
\fbox{\parbox{0.85\textwidth}{\centering\vspace{4cm}
\textbf{[그림 9.6]} 임계 형성 간 거리 $d_{\min}^*$의 비교.\\
실선: SCC (닫힘 포함). 점선: 순수 Allen-Cahn (닫힘 없음).\\
닫힘은 $d_{\min}^*$를 약 30\% 감소시켜 다중 형성의 공존 영역을 확장.
\vspace{0.5cm}}}
\caption{닫힘이 임계 형성 간 거리를 약 30\% 감소시킨다. $\beta = 30$에서 $d \approx 7$(Allen-Cahn) → $d \approx 5$(SCC).}
\label{fig:d-min-closure}
\end{figure}


\section{T-Persist-K-Unified: 통합 다중 형성 지속성 정리}
\label{sec:persist-k-unified}

이 장의 내용을 하나의 매개변수적 정리로 종합한다.

\begin{theorem}[T-Persist-K-Unified]
\label{thm:persist-k-unified}
$K$개 형성 $(u^1, \ldots, u^K)$가 다음 5개 조건을 만족하면, 모든 형성이 시간적으로 지속한다:
\begin{enumerate}
    \item[\textbf{PS}] \textbf{형성별 안정성}: 각 $u^k$가 T-Persist-1의 조건 (a)--(e)를 만족
    \item[\textbf{ND}] \textbf{비퇴화}: 결합 헤시안 $H_K$가 비퇴화 (최소 고유값 > 0)
    \item[\textbf{BC'-K}] \textbf{분지 포함}: 각 형성의 그래디언트 흐름 분지가 유지됨
    \item[\textbf{TC-K}] \textbf{전달 집중}: $K$개의 형성별 전달 계획 $M^k$가 교차 오염 없이 집중됨
    \item[\textbf{SR-$\Lambda$}] \textbf{체제 적합 조건}: $\Lambda_{\mathrm{coupling}}$의 체제에 따른 추가 조건:
    \begin{itemize}
        \item 잘 분리된 체제 ($\Lambda < 0.01$): 추가 조건 없음
        \item 약하게 상호작용 ($0.01 \leq \Lambda < 1/(K-1)$): 격리 조건(isolation condition)
        \item 강하게 상호작용 ($\Lambda \geq 1/(K-1)$): 장벽 안정성 조건
    \end{itemize}
\end{enumerate}
\end{theorem}

이 정리는 세 가지 기존 정리---T-Persist-K-Sep(잘 분리된 체제), T-Persist-K-Weak(약하게 상호작용 체제)---를 하나의 매개변수적 틀로 통합한다. $\Lambda_{\mathrm{coupling}}$의 값에 따라 적절한 조건이 자동으로 선택되며, exp46--47의 100\% 실험 검증이 이 통합의 경험적 기반을 제공한다.

\subsection{다섯 조건의 상세 해설}

각 조건의 물리적 의미와 검증 방법을 상세히 살펴보자.

\textbf{PS (형성별 안정성).} 이 조건은 각 형성이 \emph{개별적으로} 지속 가능해야 함을 요구한다. 제10장에서 다룰 T-Persist-1의 다섯 조건---지지, 분지 반경, 깊은 핵심, 약한 체제, 전달 속박---이 각 $u^k$에 대해 만족되어야 한다. 직관적으로, 다중 형성의 지속성은 ``개별 형성이 각각 살아남을 수 있을 때''만 보장된다.

\textbf{ND (비퇴화).} 결합 헤시안 $H_K$의 최소 고유값이 양수여야 한다. 이는 에너지 곡면이 극소점에서 ``그릇 모양(bowl-shaped)''임을 보장한다. 퇴화가 발생하면 극소점이 안장점으로 전환되어 안정성을 잃는다. 수치적으로, $\sigma_{\min}(H_K) > 0$은 역 반복(inverse iteration)으로 검증한다.

\textbf{BC'-K (분지 포함).} 시간적 섭동 후에도 각 형성이 원래의 극소점의 흡인 분지(basin of attraction) 내에 머무는 것을 요구한다. 이는 PS의 분지 반경 조건(b)의 K-장 버전이다. K-장에서는 교차 효과에 의해 분지 형태가 단일 장과 달라질 수 있으며, 이 변형을 고려해야 한다.

\textbf{TC-K (전달 집중).} 각 형성의 전달 계획 $M^k_{t \to s}$가 ``교차 오염'' 없이 핵심-핵심으로 집중됨을 요구한다. 구체적으로:
\begin{equation}
    \frac{\sum_{y \in \mathrm{Core}(u^j_s)} M^k_{t \to s}(x, y)}{\sum_y M^k_{t \to s}(x, y)} < \varepsilon_{\mathrm{cross}}, \quad \forall x \in \mathrm{Core}(u^k_t),\; j \neq k
    \label{eq:cross-contamination}
\end{equation}
이는 형성 $k$의 핵심에서 출발한 질량이 형성 $j$의 핵심으로 ``누출''되지 않음을 보장한다.

\textbf{SR-$\Lambda$ (체제 적합 조건).} 이 조건은 $\Lambda_{\mathrm{coupling}}$의 값에 따라 다른 강도의 요구를 부과한다. 잘 분리된 체제에서는 추가 조건이 없다---형성들이 서로를 ``보지 못하므로'' 상호작용의 효과를 고려할 필요가 없다. 약하게 상호작용하는 체제에서는 경계 대역의 겹침이 전달 집중을 방해하지 않는다는 격리 조건이 필요하다.

\subsection{통합 정리의 실험적 검증}

T-Persist-K-Unified는 exp46--47에서 체계적으로 검증되었다. 검증 프로토콜:

\begin{enumerate}
    \item 다양한 $(\lambda_{\mathrm{rep}}, d_{\min}, \beta)$ 조합으로 69개 K-장 구성 생성
    \item 각 구성에 대해 $\Lambda_{\mathrm{coupling}}$ 계산 및 체제 분류
    \item 시간적 전달 실험: 각 구성을 5시간 단계 동안 진화시키고 지속성 측정
    \item 다섯 조건(PS, ND, BC'-K, TC-K, SR-$\Lambda$) 충족 여부와 실제 지속성 일치 확인
\end{enumerate}

결과: 다섯 조건을 모두 충족하는 52개 구성에서 100\%가 지속하였고, 하나 이상의 조건을 위반하는 17개 구성 중 14개에서 지속 실패가 관찰되었다 (위반 ↔ 실패 상관계수 0.82).

% [Figure placeholder]
\begin{figure}[t]
\centering
\fbox{\parbox{0.85\textwidth}{\centering\vspace{5cm}
\textbf{[그림 9.7]} T-Persist-K-Unified의 실험적 검증.\\
가로축: $\Lambda_{\mathrm{coupling}}$. 세로축: $\mathsf{Persist}$ 점수.\\
녹색: 5조건 모두 충족 (지속 성공). 적색: 1개 이상 위반 (지속 실패 포함).\\
점선: 체제 경계 ($\Lambda = 0.01$, $\Lambda = 1/(K-1)$).
\vspace{0.5cm}}}
\caption{69개 구성에 대한 T-Persist-K-Unified 검증. 5조건 충족 시 100\% 지속, 위반 시 82\% 실패.}
\label{fig:persist-k-validation}
\end{figure}


\section{요약 및 전망}
\label{sec:multi-summary}

이 장에서 확인한 핵심 사항을 정리하자:

\begin{itemize}
    \item \textbf{단일 장의 한계}: 전역 에너지 최소점은 항상 $K = 1$이다. 다중 형성은 열역학적이 아닌 \emph{운동학적} 현상이다.

    \item \textbf{K-장 구조}: $K$개의 결합된 응집 장 + 반발 + 단체 제약이 다중 형성을 구성적으로 보장한다.

    \item \textbf{체제 분류}: 단일 매개변수 $\Lambda_{\mathrm{coupling}}$이 세 체제를 결정한다 (잘 분리된, 약하게 상호작용, 강하게 상호작용).

    \item \textbf{운동학적 기둥}: 핵생성, 준안정성($\beta^{0.89}$ 장벽), 조대화가 다중 형성의 생애를 결정한다.

    \item \textbf{닫힘의 역할}: 닫힘은 준안정 장벽을 강화하고 임계 거리를 30\% 감소시킨다.

    \item \textbf{통합 정리}: T-Persist-K-Unified가 모든 체제를 하나의 매개변수적 틀로 통합한다.
\end{itemize}

다음 장에서는 시간적 전달---형성이 시간을 가로질러 \emph{지속}하는 메커니즘---을 상세히 다룬다. 이를 통해 ``시간적 정체성은 추적 ID가 아닌 구조적 계승이다''라는 SCC의 핵심 주장이 수학적으로 어떻게 구현되는지를 볼 것이다.
